% Template:     Reporte LaTeX
% Documento:    diffThreshold sensitivity sweep
% Codificación: UTF-8

\documentclass[
	english,
	letterpaper, oneside
]{article}

\def\documenttitle {\texttt{diffThreshold} Sensitivity Sweep of the Post-Fix Hybrid Mandelbrot Renderer}
\def\documentsubtitle {Where Is the Sweet Spot, and Why Cannot the Corner Heuristic Reach It?}
\def\documentsubject {Adaptive Subdivision, Load Balancing}
\def\documentdate {\today}

\def\documentauthor {Matias Saavedra}
\def\coursename {Multicore and GPU Programming}
\def\coursecode {CISC 701}

\def\universityname {Universidad de Chile}
\def\universityfaculty {Facultad de Ciencias Físicas y Matemáticas}
\def\universitydepartment {Departamento de Ciencias de la Computación}
\def\universitydepartmentimage {departamentos/dcc}
\def\universitylocation {Santiago de Chile}

../profiling_report/template.tex

\begin{document}

\templatePagecfg
\templateFinalcfg

\inserttitle

% -----------------------------------------------------------------------
\begin{abstract}
This report sweeps the \texttt{diffThreshold} parameter of
\texttt{mandelHybrid} across the values $\{0.05, 0.1, 0.2, 0.3, 0.5\}$
on the post-fix binary, measuring how the four-corner uniformity
heuristic's only tunable knob affects end-to-end wall time, total work
generated, and the worst-case region. Three configurations are
exercised --- \texttt{CPU12} (CPU-only baseline), \texttt{GPU} (single
GPU thread), and \texttt{GPU+11CPU} (best hybrid) --- with three reps at
each \texttt{diffT}. The headline result is that the hybrid sweet spot
sits at \texttt{diffT}\,=\,$0.1$ at 52.80~s mean, a 3.3\% improvement
over the 0.5 baseline (54.61~s), with a non-monotonic tail: dropping
further to \texttt{diffT}\,=\,$0.05$ does \emph{worse} (53.45~s with
much higher variance, $\sigma{=}1.26$~s vs.\ $0.50$~s at 0.1). Pure
\texttt{GPU} improves monotonically (-0.5\% across the full range) and
\texttt{CPU12} is essentially invariant. The most surprising finding,
revealed by a \texttt{quiet=0} rerun at each \texttt{diffT}, is that
\textbf{the worst-case CPU region barely moves at all}: it remains a
single $960\times540$ region taking 14.0--15.5~s regardless of
\texttt{diffT}. The heuristic cannot subdivide this region at any
threshold because all four corners lie inside the Mandelbrot set
(\texttt{iter}=\texttt{MAXITER}), making the corner-spread identically
zero. The improvement from tightening \texttt{diffT} is therefore not
coming from shrinking the binding outlier --- it is coming from
better load balance among the other 5{,}300+ leaves.
\end{abstract}

% -----------------------------------------------------------------------
\section{Setup}

The four-corner subdivision rule used by \texttt{MandelRegion::examine()}
is:

\begin{sourcecode}{cpp}{mandelregion.cpp -- the decision rule under study.}
// either compute the pixels or break the region in 4 pieces
if (maxIter - minIter < diffThresh * maxIter
    || pixelsX * pixelsY < pixelSizeThresh)
  {
    compute (onGPU);
    ownerFrame->regionComplete ();
  }
else
  {
    // ... four-way split ...
  }
\end{sourcecode}

\noindent
At \texttt{diffThresh}~$\to$~$1$ the rule almost always passes
(everything looks uniform; only the pixel-size floor causes splits).
At \texttt{diffThresh}~$\to$~$0$ the rule almost never passes (only
identical-corner regions skip splitting; the pixel floor saturates the
behaviour). The interesting regime is in between.

\subsection{Configurations}

\begin{itemize}
    \item \textbf{Hardware and workload}: identical to the previous two
    reports --- i7-9750H + GTX~1660~Ti~Max-Q, 100 frames at
    $1920\times1080$ from the standard \texttt{spec.in}, post-fix
    binary on branch \texttt{examine\_minmax\_bugfix}, default
    \texttt{pixelSizeThresh}\,=\,$32{,}768$.
    \item \textbf{Threshold sweep}: $\{0.05, 0.1, 0.2, 0.3, 0.5\}$.
    The value $0.5$ is the legacy default; it was re-used from the
    \texttt{sweep\_results\_postfix} run rather than re-measured, to
    save 9 redundant runs.
    \item \textbf{Configurations}: three labels per \texttt{diffT} ---
    \texttt{CPU12} (\texttt{numThreads}\,=\,$12$, \texttt{gpuEnable}\,=\,$0$),
    \texttt{GPU} (\texttt{numThreads}\,=\,$1$, \texttt{gpuEnable}\,=\,$1$),
    and \texttt{GPU+11CPU} (\texttt{numThreads}\,=\,$12$, \texttt{gpuEnable}\,=\,$1$).
    The pure-GPU and best-hybrid endpoints bracket the workload, and
    the CPU-only baseline reports whether the effect is GPU-specific
    or general.
    \item \textbf{Reps and timing}: three reps per \texttt{(config, diffT)}
    pair, measured with the same end-to-end \texttt{clock\_gettime}
    bracket as the Fig~11.15 sweep ($\texttt{[total\_elapsed\_s]}$ line
    in the stderr log).
\end{itemize}

\noindent
A second, smaller pass added one \texttt{quiet=0} run per
\texttt{diffT} of \texttt{GPU+11CPU} so that the per-region
\texttt{[CPU region NNNN] size WxH compute Z.ZZZ ms} lines could be
parsed for max-region statistics. This added 5 runs.

% -----------------------------------------------------------------------
\section{Results}

\subsection{Wall Time vs.\ \texttt{diffT}}

\begin{table}[H]
    \centering
    \caption{Mean wall-clock time per configuration and \texttt{diffT}
             ($n=3$ each). Best per config in bold.}
    \label{tab:wall}
    \begin{tabular}{lrrrrr}
        \hline
        \textbf{Config} & \textbf{0.05} & \textbf{0.10} & \textbf{0.20} & \textbf{0.30} & \textbf{0.50} \\
        \hline
        CPU12     & \textbf{161.47} & 161.96 & 162.59 & 162.24 & 162.76 \\
        GPU       & \textbf{77.60} & 77.81 & 77.82 & 77.83 & 78.03 \\
        GPU+11CPU & 53.45 & \textbf{52.80} & 52.90 & 53.21 & 54.61 \\
        \hline
    \end{tabular}
\end{table}

\begin{table}[H]
    \centering
    \caption{Population standard deviations per configuration and
             \texttt{diffT}, in seconds.}
    \label{tab:std}
    \begin{tabular}{lrrrrr}
        \hline
        \textbf{Config} & \textbf{0.05} & \textbf{0.10} & \textbf{0.20} & \textbf{0.30} & \textbf{0.50} \\
        \hline
        CPU12     & 0.47 & 0.71 & 0.32 & 0.94 & 2.44 \\
        GPU       & 0.03 & 0.08 & 0.08 & 0.09 & 0.21 \\
        GPU+11CPU & 1.26 & 0.50 & 0.25 & 0.75 & 0.39 \\
        \hline
    \end{tabular}
\end{table}

\insertimage{../difft_sweep/difft_sweep.png}{width=\linewidth}
            {Left: mean wall time per config across \texttt{diffT} (log
             x-axis), error bars showing standard deviation. Right:
             zoom on the \texttt{GPU+11CPU} hybrid with the
             worst-case CPU region time overlaid on a secondary
             axis. \label{fig:sweep}}

\subsection{Work Generated vs.\ \texttt{diffT}}

\begin{table}[H]
    \centering
    \caption{Total leaf regions generated per run, averaged across 3
             reps. The total is workload-independent and bounded above
             by the pixel-size floor.}
    \label{tab:leaves}
    \begin{tabular}{lrrrrr}
        \hline
        \textbf{Config} & \textbf{0.05} & \textbf{0.10} & \textbf{0.20} & \textbf{0.30} & \textbf{0.50} \\
        \hline
        \texttt{CPU12} (CPU leaves)        & 5{,}323 & 5{,}323 & 5{,}320 & 5{,}290 & 5{,}152 \\
        \texttt{GPU} (GPU leaves)          & 5{,}323 & 5{,}323 & 5{,}320 & 5{,}290 & 5{,}152 \\
        \texttt{GPU+11CPU} CPU/GPU split   & 1381/3942 & 1334/3989 & 1334/3986 & 1408/3882 & 1339/3813 \\
        \texttt{GPU+11CPU} total           & 5{,}323 & 5{,}323 & 5{,}320 & 5{,}290 & 5{,}152 \\
        \hline
    \end{tabular}
\end{table}

\subsection{Worst-Case Region vs.\ \texttt{diffT}}

The single \texttt{quiet=0} rep per \texttt{diffT} of \texttt{GPU+11CPU}
records the maximum \texttt{[CPU region]} compute time and the pixel
dimensions of that region:

\begin{table}[H]
    \centering
    \caption{Worst-case CPU region per \texttt{diffT}, single \texttt{quiet=0}
             rep, \texttt{GPU+11CPU} configuration. Single-rep wall
             times are noisier than Table~\ref{tab:wall}.}
    \label{tab:worst}
    \begin{tabular}{rrlrr}
        \hline
        \textbf{diffT} & \textbf{Wall (s)} & \textbf{Worst CPU region} & \textbf{Worst kernel (ms)} & \textbf{Worst kernel size} \\
        \hline
        0.05 & 51.15 & 14.00~s @ $960\times540$ & 137.9 & $480\times270$ \\
        0.10 & 53.08 & 15.47~s @ $960\times540$ & 139.8 & $480\times270$ \\
        0.20 & 54.00 & 15.35~s @ $960\times540$ & 140.7 & $480\times270$ \\
        0.30 & 54.75 & 15.41~s @ $960\times540$ & 140.7 & $480\times270$ \\
        0.50 & 54.34 & 15.29~s @ $960\times540$ & 141.1 & $480\times270$ \\
        \hline
    \end{tabular}
\end{table}

% -----------------------------------------------------------------------
\section{Analysis}

\subsection{The Hybrid Sweet Spot Is \texttt{diffT}\,=\,$0.1$}

The clearest finding is that the best hybrid configuration is not at
the legacy default (0.5) nor at the most aggressive value (0.05) but at
\texttt{diffT}\,=\,$0.1$, with a mean of 52.80~s. Relative to the 0.5
baseline:

\begin{itemize}
    \item $0.5 \to 0.3$: $-1.40$~s ($-2.6\%$).
    \item $0.5 \to 0.2$: $-1.71$~s ($-3.1\%$).
    \item $0.5 \to 0.1$: $-1.81$~s ($-3.3\%$).
    \item $0.5 \to 0.05$: $-1.16$~s ($-2.1\%$), but with $\sigma=1.26$~s
    vs.\ $0.39$~s at 0.5. The mean improved but the worst rep ran
    \emph{worse} than the 0.5 baseline.
\end{itemize}

\noindent
The point of inflection between 0.1 and 0.05 is the only non-monotonic
piece of the curve. Past 0.1 the threshold no longer produces extra
useful splits (the pixel-size floor takes over: total leaves saturate
at exactly 5{,}323 by \texttt{diffT}\,=\,$0.1$) but it does increase
queue-scheduling jitter, which shows up in the variance.

\subsection{Pure GPU Benefits Monotonically (But Tinily)}

The \texttt{GPU} configuration improves from 78.03~s at 0.5 to 77.60~s
at 0.05 --- a 0.43~s ($-0.5\%$) improvement that is small in fraction
but real in $\sigma$ units: the standard deviation at every
\texttt{diffT} for \texttt{GPU} is below 0.1~s, so the 0.43~s gap is
$>5\sigma$ of either endpoint. The mechanism is the same as on the
hybrid side --- smaller, more uniform regions mean less time spent on
the worst-case kernel launch --- but with only one GPU thread there is
no load-balance benefit to capture, only the modest reduction in tail
kernel duration (from 141.1~ms to 137.9~ms; see
Table~\ref{tab:worst}).

\subsection{CPU-Only Is Almost Invariant}

\texttt{CPU12} runs in 161.5--162.8~s across the entire \texttt{diffT}
range, with most of the variation absorbed by run-to-run noise. The
mechanism: with 12 CPU threads sharing roughly 5{,}300 leaves of work
totalling $\sim$2000~CPU-thread-seconds, the per-leaf granularity is
already fine enough at any \texttt{diffT} for the dynamic queue to
balance well. The bottleneck in CPU-only mode is total
work-divided-by-threads, not load imbalance, so the small +3\% leaves
created by going from \texttt{diffT}\,=\,$0.5$ to $0.05$ adds 0.5\%
to the total compute time and saves a similar amount on the tail --- a
wash.

\subsection{The Surprise: The Outlier Is Unbreakable}

The worst single CPU region remains a $960 \times 540$ patch
\textbf{at every \texttt{diffT} tested}, taking 14.0--15.5~s. Five
points of \texttt{diffT} from 0.05 to 0.5 (a 10$\times$ range) move
its compute time by less than 1.5~s. This is the central surprise of
the experiment, and it is also the explanation for why all the
preceding wall-time improvements are so small:

\textbf{The corner-spread heuristic cannot see this region as a split
candidate at any \texttt{diffT}.} A $960 \times 540$ leaf is a quarter
of an early-zoom frame; its four corners almost certainly all lie
inside the Mandelbrot set, where \texttt{diverge()} returns the cap
\texttt{MAXITER}\,=\,$10{,}000$. With identical corner values, the
spread $\texttt{maxIter} - \texttt{minIter} = 0$, and the decision
rule \texttt{0 < diffThresh * maxIter} is true for every positive
\texttt{diffThresh}. The region is classified as uniform and dispatched
to \texttt{compute()} in one block, where it sits on a single CPU
thread for $\sim$15~s while the other threads complete the remaining
$\sim$5{,}320 leaves and idle.

Lowering \texttt{diffT} does not help because the corners are
\emph{identical}, not just close. The diffT knob controls how much
\emph{near}-uniformity is acceptable; it cannot reject \emph{exact}
uniformity. The only way to subdivide this region with the existing
heuristic is to lower \texttt{pixelSizeThresh} far enough that
$960 \times 540 = 518{,}400$ pixels falls below the floor. That would
force-split the region geometrically, which would attack the outlier
but would also force-split the much larger interior of the set (a
substantial fraction of every frame), generating overhead that very
likely overwhelms the gain.

\subsection{So Where Does the 3.3\,\% Hybrid Improvement Come From?}

Not from shrinking the worst region, since it is essentially
unchanged. The improvement comes from the other 5{,}320 leaves.
Specifically, at \texttt{diffT}\,=\,$0.5$ the renderer produces 5{,}152
leaves, but the bigger averaged-size of those leaves means the GPU
thread spends slightly more time per kernel launch and the CPU threads
spend slightly more time per region. At \texttt{diffT}\,=\,$0.1$ the
renderer produces 5{,}323 leaves --- 171 more, $+3.3\%$ --- each
\emph{slightly} smaller on average. Three effects compound:

\begin{enumerate}
    \item \textbf{More uniform region sizes} mean the tail of the leaf
    distribution (the leaves the CPU threads chew on) is closer to
    the median, so any given CPU thread finishes its last leaf closer
    to the slowest thread's completion time. Per-thread wall times
    tighten.
    \item \textbf{The GPU consumes leaves $\sim$40$\times$ faster than
    the CPU}, so the GPU thread is always returning to the queue
    looking for more work. More leaves at finer granularity means the
    GPU stays usefully busy for longer; its idle tail (the time it
    spends waiting for the last CPU thread) shrinks.
    \item \textbf{The worst-case region's $\sim$15~s wall is masked by
    other workers}. When that region runs on, say, thread~7, the
    other 10 CPU threads continue chewing on the queue. The wall time
    becomes \texttt{max(15.3~s,~time-to-drain-remaining-queue)}. More
    leaves means the queue takes slightly longer to drain, so the
    15.3~s is more often \emph{below} the queue-drain time and ceases
    to be the binding constraint for that rep.
\end{enumerate}

\noindent
The third effect explains why \texttt{diffT}\,=\,$0.05$ has high
variance: the queue-drain time becomes very close to the worst-region
time, and small scheduling differences flip which side of that
inequality dominates the rep. Sometimes the queue takes longer
($\Rightarrow$ stable, fast rep); sometimes it doesn't
($\Rightarrow$ the 15~s sets the wall and the rep is slow).

\subsection{Why The Wall Improvement Plateaus}

Five points of \texttt{diffT} change the wall by 1.81~s. The same
order-of-magnitude as the noise in a single rep at the most aggressive
threshold. The remaining 52--53~s of wall time is the irreducible
constant on this platform with this binary: $\sim$15~s for the worst
region plus $\sim$37~s of remaining-queue drain. To improve further
the renderer needs:

\begin{itemize}
    \item \textbf{To split the worst region}. The corner heuristic
    will never do this. Required: edge-midpoint sampling (the centre
    of the worst $960 \times 540$ region likely has a non-trivial
    iteration value; sampling it would detect interior non-uniformity)
    or a geometric/pixel-floor-only force-split for any region whose
    corners are all \texttt{MAXITER}.
    \item \textbf{To use cardioid and bulb interior certificates.} For
    truly-interior regions, this skips iteration entirely. The
    payoff in worst-case savings is bounded by however much of the
    outlier region is genuinely inside vs.\ just appearing to be from
    the corners.
    \item \textbf{Asynchronous CUDA streams}. Independent of the
    heuristic, recovers up to $\sim$13~ms of per-region GPU idle.
\end{itemize}

% -----------------------------------------------------------------------
\section{Recommendations}

The cheapest configuration change available today is to lower the
default \texttt{diffThreshold} from $0.5$ to $0.1$:

\begin{itemize}
    \item Best-hybrid wall: $54.61 \to 52.80$~s ($-3.3\%$).
    \item Pure-GPU wall: $78.03 \to 77.81$~s ($-0.3\%$, but it is
    still GPU's best).
    \item CPU-only wall: unchanged within noise.
    \item Run-to-run variance for the hybrid actually \emph{tightens}
    at $0.1$ versus $0.5$ ($\sigma = 0.50$~s vs.\ $0.39$~s --- nearly
    equal --- but the $\sigma = 2.44$~s outlier observed on
    \texttt{CPU12} at $0.5$ does not recur at $0.1$).
\end{itemize}

\noindent
Avoid $\texttt{diffT}\,=\,$0.05$, which trades the small mean improvement
for $2.5 \times$ higher run-to-run variance --- predictable performance
is more valuable than a marginal mean.

The diffT knob is now functionally exhausted on this workload. To
shrink the wall time further the heuristic itself must be replaced or
augmented. The three modifications listed at the end of
\texttt{postfix\_report.tex} (edge-midpoint sampling, cardioid
certificates, async streams) remain the right candidates; this report
now provides quantitative evidence for the urgency of the first one,
since the 15~s outlier is demonstrably immune to any tightening of the
existing rule.

% -----------------------------------------------------------------------
\section{Conclusion}

The \texttt{diffThreshold} parameter has a real but modest effect on
end-to-end wall time. For the best hybrid configuration on this
machine, the sweet spot is \texttt{diffT}\,=\,$0.1$ at 52.80~s, a
3.3\,\% improvement over the legacy default. The mechanism is better
load balance among the non-outlier leaves, not elimination of any
worst-case region.

The surprise of the experiment is that the worst-case CPU region ---
a single $960 \times 540$ patch consuming $\sim$15~s on one CPU
thread --- is invariant across the entire \texttt{diffT} sweep. Its
four corners are all at \texttt{MAXITER}, making the corner-spread
identically zero and the uniformity rule pass at every threshold.
This region is the true binding constraint on the wall time, and no
amount of \texttt{diffT} tuning can subdivide it. The four-corner
heuristic has reached its ceiling on this workload; meaningful future
gains require either sampling more points inside each region (so that
interior filaments cannot hide behind agreeing corners) or analytic
interior tests (so that genuinely interior regions are recognised and
constant-filled rather than iterated).

\end{document}
